\documentclass[11pt]{article}
\usepackage[review]{acl}
\usepackage{times}
\usepackage{latexsym}
\usepackage[T1]{fontenc}
\usepackage[utf8]{inputenc}
\usepackage{microtype}
\usepackage{inconsolata}
\usepackage{graphicx}
\usepackage{stfloats}
\usepackage{xurl}
\graphicspath{{figures/emnlp/}{./}}
\usepackage{booktabs}
\usepackage{amsmath}

% ===========================================================================
% SCAFFOLD v2 -- final numbers + GPT round-2 language discipline baked in.
% Page budget in comments (8pp body; refs unlimited; Limitations REQUIRED + not
% counted; appendix after refs, not counted). Figures: figures/emnlp/ (run
% scripts/emnlp_generate_figures.py on the pod). CALIBRATED CONFIDENCE: strong
% where evidence is strong (flip is real, position+stage-specific, survives
% held-out name split; dip in continuous LD; controlled dip predicts), hedged on
% inference (detection/suppression LABELS; composition not shown; from-scratch =
% dip not flip). Verbs: "evidence supports", "explains/predicts", "consistent
% with" -- NOT "requires"/"proves"/"the mechanism".
% ===========================================================================

\title{Detection Before Suppression: Sequential Circuit Formation\\
Explains a Below-Chance Dip in Language Model Training}

\author{First Author \\ Affiliation \\ \texttt{email} \\\And
        Second Author \\ Affiliation \\ \texttt{email} \\}

\begin{document}
\maketitle

% ============================ ABSTRACT =====================================
% Plain-language, NLP-reviewer-readable in ~20s. No "residual stream"/"ablation".
\begin{abstract}
During training, language models can temporarily perform \emph{below chance} on
structured forced-choice contrasts before recovering, a regime invisible to the loss curve.
Studying indirect-object identification in Pythia models at three scales, we
find this dip is not a metric artifact: both accuracy and the continuous logit
difference between the correct and competing answers fall below their neutral
values over a contiguous interval. Our evidence supports an explanation in which
two components come online \emph{sequentially}: an early duplicate-detection
signal and a later component with a duplicate-suppressing causal effect.
During the gap, detection signals are present while the suppression effect is
absent, so the model favors the repeated (wrong) name. The signature is a
reversal: the same intervention that \emph{improves} the model's behavior during
the dip \emph{harms} it after recovery. The reversal is position- and
stage-specific, survives held-out names, replicates on greater-than, and is
absent in a subject--verb agreement boundary case, consistent with the shortcut
not appearing to be reused there. Finally, we predictably modulate the dip
by controlling the training signal. A component's causal role can
thus reverse during training, so interpreting only mature checkpoints can
misrepresent how linguistic behaviors form.
\end{abstract}

% ============================ 1. INTRODUCTION (~1.0 pp) ====================
\section{Introduction}
% Hook with the dip (Fig 1). TWO linked claims: (1) PHENOMENON -- below-chance
% regime, in accuracy AND continuous LD; (2) THESIS -- a component whose causal
% role REVERSES explains it (detection before suppression).
% EMNLP-RELEVANCE (mandatory sentence): "interpreting a mature language model
% alone can miss transient but causally important stages in how linguistic
% circuits form." NLP-translation of tasks: IOI = coreference-like name tracking;
% greater-than = numeric comparison expressed in text; SVA = canonical syntactic
% agreement. Contributions (4-5 bullets). Frame the delta over any prior
% single-scale observation WITHOUT naming a prior paper (ARR anonymity): here =
% 3 scales, 2 tasks causally, sequential-formation mechanism with matched causal
% rigor, held-out-name + selection-bias controls, and predictable control.
Mechanistic interpretability usually studies a model at a single, mature
checkpoint. But the circuits that implement a behavior are not present at
initialization; they form over training, and the path matters. We document a
case where it matters sharply: during training, language models pass through a
window in which they perform \emph{below chance} on a structured task, actively
preferring the wrong answer, before recovering to near-perfect accuracy. The
window is invisible to the loss curve, which decreases monotonically throughout.

We study this on indirect-object identification (IOI), a controlled
name-tracking task, in Pythia models at three scales
(Figure~\ref{fig:dip_signflip}). The dip is not an artifact of a binarized
metric: the continuous logit difference between the correct and competing
answers also falls below zero over a contiguous interval. We then give causal
evidence for why it occurs. Our evidence supports a two-component explanation: an
early duplicate-detection signal and a later component with a duplicate-suppressing causal effect, the two
coming online at different times. During
the gap between them, duplicate-detection signals are present before the
later suppression effect becomes active, so the model prefers the duplicate. The
signature is a reversal: the same intervention on the duplicate's representation
that \emph{improves} the model's output during the dip \emph{harms} it after
recovery.

This yields two linked claims. The phenomenon is a below-chance dip, visible in
both accuracy and the continuous logit difference. The explanation is a reversal
in the causal role of an implicated component across training: a component that
looks harmful mid-training supports correct behavior at maturity. The methodological
consequence is direct: interpreting a model at a single mature checkpoint can
misrepresent how a linguistic behavior forms, because a component's role is not
fixed over training.

The behaviors we study are controlled linguistic behaviors rather than arbitrary probes. Our three tasks
instantiate controlled linguistic phenomena: coreference-like name tracking
(IOI), numeric comparison expressed in text (greater-than), and syntactic
agreement under an attractor (subject--verb agreement). The dip and
reversal recur across them in a predictable, falsifiable way. We make the
following contributions.

\begin{itemize}\setlength{\itemsep}{1pt}
\item We document a below-chance dip during training on IOI at three model
scales, in both accuracy and the continuous logit difference, which argues
against a metric-artifact explanation.
\item We show that the causal role of the duplicate-name representation reverses
across the dip (helpful when perturbed during the dip, harmful at maturity), and
that the reversal is position- and stage-specific, survives a held-out name
split, and survives confidence intervals clustered by template and name pair.
\item We trace the dip to the sequential formation of a duplicate-detection
signal before a suppression effect, with causal measurements on both sides.
\item We show the reversal generalizes causally to greater-than and is absent in
a subject--verb agreement boundary case, instantiating a falsifiable prediction.
\item We predictably modulate the behavioral dip in independently trained
from-scratch models by varying the training signal.
\end{itemize}

\begin{figure*}[!t]
  \centering
  \includegraphics[width=\textwidth]{fig1_dip_and_signflip.pdf}
  \caption{\textbf{The same intervention reverses sign across training.}
  \textbf{(a)} IOI accuracy dips below chance (shaded window) at all three scales.
  \textbf{(b)} The dip is not a metric artifact: the continuous logit difference
  also goes negative. \textbf{(c)} The S2 intervention \emph{improves} behavior
  during the dip ($\Delta$LD $>$ 0) and \emph{harms} it after recovery
  ($\Delta$LD $<$ 0), so the causal effect of perturbing the affected representation
  reverses over training.}
  \label{fig:dip_signflip}
\end{figure*}

% ============================ 2. SETUP (~0.7 pp) ===========================
\section{Setup}
\label{sec:setup}
\paragraph{Task.}
We study indirect-object identification (IOI), the task isolated by
\citet{wang2023ioi} as a tractable window into how transformers track entities.
Given a prompt such as ``When Mary and John went to the store, John gave a drink
to,'' a competent model should continue with the indirect object \emph{Mary}
(IO) rather than the repeated subject \emph{John} (S). Each prompt contains three
name positions: the indirect object IO (\emph{Mary}), the first mention of the
subject (S1), and its repetition (S2), the two occurrences of \emph{John}. IOI
is coreference-like in that the model must use sentence structure, specifically
that one name has already appeared twice, to decide which entity is the novel
referent that fills the final slot. We instantiate a curated set of ten templates
with single-token names, which keeps scoring unambiguous, since a single
next-token prediction decides each prompt, and yields 300 prompts per checkpoint
(30 per template). We analyze the released checkpoints of Pythia-160M, 410M, and
1B \citep{biderman2023pythia}, which expose a model's full training trajectory
rather than only its final state.

\paragraph{Metrics.}
We measure two quantities at the final token position. \emph{Accuracy} is the
fraction of prompts on which the model assigns higher probability to IO than to
S, a forced choice between the two candidate names. The \emph{logit difference}
$\mathrm{LD}=\mathrm{logit}(\text{IO})-\mathrm{logit}(\text{S})$ is the
corresponding continuous quantity, positive when the model favors the correct
answer and negative when it favors the repeated name. We rely on the logit
difference because it is continuous and signed: it separates a model that is
merely uncertain (LD near zero) from one that \emph{actively} prefers the wrong
answer (LD negative), and it avoids the thresholding artifacts by which a
discontinuous metric can appear to change sharply when the underlying behavior
does not \citep{schaeffer2023mirage}. \emph{Below chance} refers to this
IO-versus-S forced-choice contrast: the model places higher probability on S
than on IO, equivalently below 50\% accuracy on the binary choice, not to
below-random prediction over the open vocabulary, for which 50\% is not the
relevant baseline.

\paragraph{Causal intervention.}
To measure the causal role of the repeated name we intervene on the model's
internal computation. At each layer a transformer maintains a \emph{residual
stream}, a per-token vector that accumulates the contributions of the attention
and feed-forward sublayers and carries information forward. Our probe replaces
the residual stream at the S2 token position with the residual from a matched
\emph{control} prompt that is identical except that S2 is a third, unrelated
name, and measures the resulting change in logit difference,
$\Delta\mathrm{LD}=\mathrm{LD}_{\text{patched}}-\mathrm{LD}_{\text{base}}$. This
asks what the model's answer would have been had the duplicate-name signal at S2
been replaced by a neutral one; a large $\Delta\mathrm{LD}$ means the output
depends causally on what the model computed at that position. We apply the
intervention over a contiguous window of layers matched to model depth:
$\{3,4,5\}$ for 160M, $\{6,7,8,9,10\}$ for 410M, and $\{4,5,6,7,8\}$ for 1B
(roughly 25--45\% relative depth). Layer windows are fixed before held-out
evaluation and are not chosen using the reported effects;
Appendix~\ref{sec:appendix} shows the result is stable across neighboring
windows.

\paragraph{Defining the dip, and guarding against bias.}
We define a \emph{below-chance dip} as a contiguous post-initialization interval
over which mean LD is significantly negative (its bootstrap confidence interval
excludes zero) and forced-choice accuracy is below 50\%; for each model we report
the dip onset, floor checkpoint, minimum accuracy, and minimum LD with confidence
intervals. Because a checkpoint chosen as the floor could be an outlier, we guard
against checkpoint-selection bias: we split the single-token name pool into two
disjoint halves, identify dip checkpoints using only unpatched base accuracy on
the selection half, and report all intervention effects on the held-out half,
across the whole dip interval rather than only at the floor. Unless stated
otherwise, confidence intervals are 95\% bootstrap intervals clustered by
template and by name pair, so that shared structure among prompts does not
inflate apparent precision.

% ============================ 3. THE DIP AND THE SIGN FLIP (~1.5 pp) =======
\section{A Below-Chance Dip and a Causal Sign Flip}
\label{sec:spine}
Across Pythia-160M, 410M, and 1B, IOI accuracy is non-monotonic during
training. After an initial rise it falls below the forced-choice baseline of
50\%, reaching floors of 31.7\%, 29.3\%, and 36.3\% (at step 2000 for 160M and
410M, step 1000 for 1B), then recovers to 99.7--100\% at maturity
(Figure~\ref{fig:dip_signflip}a). The same non-monotonicity appears in the
continuous logit difference $\mathrm{LD}=\mathrm{logit}(\text{IO})-\mathrm{logit}(\text{S})$:
mean LD is negative over the floor interval at every scale
(Figure~\ref{fig:dip_signflip}b), with bootstrap confidence intervals excluding
zero. The dip is therefore not an artifact of binarizing a noisy margin: the
model assigns higher probability to the repeated name S than to the correct
indirect object IO. IOI is a controlled, coreference-like name-tracking task:
the model must identify which of two named entities is the novel referent, so the
dip is a transient period in which it systematically favors the wrong name on this
controlled contrast. Table~\ref{tab:dipstats} summarizes the dip at each scale.
The phenomenon is not specific to these checkpoints or to a single random seed:
all nine PolyPythias-160M seeds dip below chance at the same training step
(floor accuracy 18--37\%; Appendix~\ref{sec:appendix-models},
Figure~\ref{fig:polyseeds}), the causal sign flip reproduces in PolyPythias
seed-1 and in an independently trained, different-architecture model (Stanford
GPT-2, floor accuracy 12\%). We also observe suggestive behavioral evidence of
the dip in OLMo-1B, a different model family, without a causal intervention.
Appendix~\ref{sec:appendix-models} reports per-model numbers.

\begin{table}[t]
  \centering\small
  \setlength{\tabcolsep}{4pt}
  \begin{tabular}{lccccc}
    \toprule
    Model & Onset & Floor & Floor & Floor & Final \\
          & (step) & (step) & acc. & LD & acc. \\
    \midrule
    160M & 512  & 2000 & 31.7\% & $-0.75$ & 99.7\% \\
    410M & 1000 & 2000 & 29.3\% & $-0.70$ & 99.7\% \\
    1B   & 1000 & 1000 & 36.3\% & $-0.54$ & 100.0\% \\
    \bottomrule
  \end{tabular}
  \caption{Dip statistics across scale: the training step at which forced-choice
  accuracy first falls below 50\% (onset), the floor checkpoint, and the floor
  accuracy, floor logit difference, and final accuracy. All three models dip
  below chance with negative LD and then recover to near-perfect accuracy.}
  \label{tab:dipstats}
\end{table}

We probe the role of the duplicate-name (S2) residual using the
S2-to-control intervention of Section~\ref{sec:setup}
(Appendix~\ref{sec:appendix} confirms robustness to neighboring layer windows).
The effect $\Delta\mathrm{LD}=\mathrm{LD}_{\text{patched}}-\mathrm{LD}_{\text{base}}$
reverses sign across training (Figure~\ref{fig:dip_signflip}c). At the floor,
perturbing the S2 representation raises LD ($\Delta\mathrm{LD}=+0.95$ at 160M); at
maturity it lowers LD sharply ($-4.12$). The reversal holds at 410M ($+0.10\to-3.63$) and
1B ($+0.43\to-3.49$), all confidence intervals excluding zero. The same S2 representation whose
perturbation harms mature performance improves performance during the dip.

On the disjoint held-out name split (Section~\ref{sec:setup}), the flip
survives: held-out $\Delta\mathrm{LD}=+1.13\ [1.04,1.23]$ at the dip and
$-4.09\ [-4.30,-3.89]$ at maturity, consistent across the dip interval and not
only at the floor.

The reversal is specific to the S2 position and to the training stage.
Table~\ref{tab:control_battery} perturbs five positions at the dip and at
maturity.

\begin{table*}[t]
  \centering
  \small
  \setlength{\tabcolsep}{10pt}
  \begin{tabular}{lccc}
    \toprule
    Position (perturbation) & $\Delta$LD (dip) & $\Delta$LD (mature) & Flip? \\
    \midrule
    \textbf{S2}, counterfactual & $\mathbf{+0.95}$ [.86,1.05] & $\mathbf{-4.12}$ [-4.28,-3.96] & \textbf{yes} \\
    \textbf{S2}, noise          & $\mathbf{+0.92}$ [.84,1.00] & $\mathbf{-4.13}$ [-4.31,-3.94] & \textbf{yes} \\
    S1, noise                   & $+2.15$ [2.06,2.24] & $+2.08$ [1.88,2.27] & no \\
    IO, noise                   & $-3.54$ [-3.73,-3.34] & $-9.06$ [-9.31,-8.82] & no \\
    structural, noise           & $-0.81$ [-.90,-.73] & $-4.45$ [-4.67,-4.23] & no \\
    \bottomrule
  \end{tabular}
  \caption{Control battery (Pythia-160M; 95\% bootstrap CIs, $n{=}300$). Only the
  S2 position reverses sign, and it reverses identically for a counterfactual
  name and norm-matched noise. The diagnostic is the stage-dependent
  \emph{reversal}, not the effect magnitude: non-S2 positions have large effects
  but none reverse.}
  \label{tab:control_battery}
\end{table*}
Only S2 reverses sign, and it reverses identically whether S2 is replaced by a
counterfactual name or by norm-matched noise. We treat the counterfactual name
as the primary in-distribution intervention and the norm-matched noise as a
nonspecific perturbation control. The diagnostic comparison is S2 versus
non-S2, not counterfactual versus noise: generic disruption is insufficient,
since the other positions produce large effects but none reverse across the
training stage. We do not interpret these large non-S2 effects as failures of
specificity; the specificity claim concerns the predicted stage-dependent sign
reversal. Perturbing S1, a second copy of the repeated name, raises LD at both
stages, consistent with weakening the repeated-name signal; perturbing IO
removes target information and is damaging at both stages; perturbing a
structural position is a nonspecific disruption that is mildly damaging at both
stages, and none reverses. The controls thus show that a perturbation can be
large without being diagnostic: what is specific to S2 is not effect size but
role reversal.

% ============================ 4. MECHANISM (~1.6 pp) =======================
\section{Detection Before Suppression}
\label{sec:mechanism}
Our analysis uses three standard interpretability primitives, which we state
plainly for readers outside the subfield. An \emph{induction head} is an
attention head that implements prefix matching: having seen the bigram
\texttt{A\,B} earlier, it attends from a later occurrence of \texttt{A} to the
token that followed it, and is the canonical mechanism by which transformers copy
a previously seen token \citep{elhage2021mathematical,olsson2022induction}; its
prevalence is summarized by an induction score. \emph{Mean-ablation} of a head
replaces its output with its average over the dataset, removing its
input-dependent contribution, so the resulting change in logit difference
measures how much the model's answer depends on that head. The \emph{logit lens}
projects a layer's residual stream directly through the unembedding, reading off
which answer each layer would predict if it were the last
\citep{nostalgebraist2020logitlens}. We next characterize candidate components
whose timing can explain the dip, tracking three quantities across training in
Pythia-160M (Figure~\ref{fig:mechanism}). First, the
prefix-matching (induction) score rises sharply early, from 0.04 to 0.94 by step
1000, and stays high throughout the below-chance interval. Second, we locate the
attention head with the largest mean-ablation effect on mature LD, namely L8H9, the
strongest of five heads whose ablation reduces the mature logit difference (the
five, by effect on mature LD, are L8H9 $-2.66$, L3H0 $-1.78$, L4H11 $-1.54$,
L6H2 $-0.78$, and L1H8 $-0.66$; L8H9 alone accounts for the largest share). We
select
this head using only mature-checkpoint behavior, then ask whether its causal
effect is already present during the dip. Tracking its effect across training,
its ablation changes LD by approximately zero during
the dip (steps 1000--2000) and grows to $-2.66$ at maturity, approximately 70\%
of the mature base LD magnitude of $+3.79$. Third, an unembedding projection of
intermediate layers favors S at every layer during the dip, with a late layer
(L8) first favoring IO only at recovery; this layer marker is clean at 160M and
noisier at 410M. Throughout, language-model loss on held-out text decreases
monotonically.

\begin{figure}[t]
  \centering
  \includegraphics[width=\columnwidth]{fig2_mechanism_timeline.pdf}
  \caption{Pythia-160M. The induction score is already high during the
  below-chance window (shaded); the dominant S-inhibition head's causal effect
  appears only at recovery; LM loss falls monotonically throughout.}
  \label{fig:mechanism}
\end{figure}

As independent evidence that the duplicate is detected during the dip, we train a
cross-validated linear probe to predict, from the final-token representation,
whether the prompt contains a repeated name (a structural property, distinct
from which name is the answer).
During the dip (steps 1000--2000) this probe already reaches 93\% accuracy on a
balanced held-out set of prompts with and without a repeated name (chance 50\%),
in middle and late layers (layers 5--11 at step~2000), rising toward ceiling by
recovery. The duplicate is therefore linearly and generalizably
represented while the model's IOI behavior is still below chance, so detection is
present before the model uses it correctly (Figure~\ref{fig:dupprobe}). This contrasts with a probe for the
answer itself (which name is the indirect object), which generalizes poorly to
held-out names (Appendix~\ref{sec:appendix-probe}), as expected if the answer
mapping is partly memorized while duplicate detection is a learned structural
feature.

\begin{figure}[t]
  \centering
  \includegraphics[width=\columnwidth]{fig6_dupprobe_by_layer.pdf}
  \caption{During the dip (Pythia-160M, step~2000), a cross-validated probe for
  whether a prompt contains a repeated name is already well above the balanced
  50\% baseline in middle and late layers (blue), confirming the duplicate is
  detected before the model uses it correctly. Error bars are cross-validation
  standard deviations.}
  \label{fig:dupprobe}
\end{figure}

We refer to the early, prefix-matching component as \emph{detection} and the
later, ablation-sensitive component as \emph{suppression}. The two are
temporally dissociated: detection signals are high while the suppression effect
is near zero, and the suppression effect emerges only as accuracy recovers. This
temporal dissociation suggests that the dip corresponds to an interval in which
the duplicate name is detected before its later suppression becomes causally
active. Informally, the model learns to notice the duplicate before it learns to
ignore it. We call L8H9 a dominant suppression head rather than the suppression
mechanism: its ablation accounts for a large fraction of the mature effect, but
causal effects in nonlinear networks are not additive, and other heads
contribute.

Our evidence establishes a temporal and causal dissociation, not a complete
mediated circuit. Duplicate-detection signals are already present during the dip,
while the strongest identified suppression component is causally inert until
recovery. We do not claim to have shown the full path by which detection is
routed into suppression; path patching or mediation analyses would be needed to
establish that composition directly. The account is therefore that sequential
formation explains the observed dip and sign flip, while the exact wiring between
the components remains an open mechanistic question. A stronger necessity test
would ablate the high-induction components during the dip and ask whether the
below-chance preference weakens; we leave this targeted test to future work.

Taken together, the results support a sequential-formation account. The dip is
real on both accuracy and the continuous LD (Section~\ref{sec:spine}). The S2
intervention reverses sign across training, on held-out names, while no other
position does. Detection-like induction signals are already high during the dip.
Suppression-like causal effects are absent during the dip and emerge only at
recovery. The same reversal appears in greater-than but not in subject--verb
agreement, where the shortcut does not appear to be reused
(Section~\ref{sec:general}). This pattern is difficult to explain as a metric
artifact, generic activation damage, or checkpoint noise.

We consider these alternatives directly. The continuous LD going negative, not
only the binarized accuracy, argues against a metric-artifact explanation. The
monotonic language-model loss argues against a loss-level double-descent
explanation. Memorization is unlikely, as the flip generalizes to held-out names,
and the behavioral dip can be reproduced in independently trained from-scratch
models (Section~\ref{sec:control}). Prompt-distribution shift is controlled by
the fixed curated protocol. Checkpoint noise is implausible given the
position-specific reversal, replicated across three scales and, for the
controlled dip, across seeds.

% ============================ 5. GENERALIZATION (~1.1 pp) ==================
\section{Generalization and a Boundary}
\label{sec:general}
The account predicts the sign flip should recur for other tasks whose early
shortcut is later repurposed. Greater-than, a numeric comparison expressed in
text, where the model completes a sentence such as ``The war lasted from 1741 to
17{\_}{\_}'' with a two-digit year exceeding the start, shows the same causal
reversal at all three scales (Figure~\ref{fig:general}a). Patching the start-year
representation toward a
control raises the correct-year margin during the dip ($\Delta\mathrm{LD}=+0.08,
+0.10, +0.32$ at 160M, 410M, 1B) and lowers it at maturity ($-0.74, -0.67,
-0.73$), all confidence intervals excluding zero. As a negative control,
next-token loss on natural text remains monotonic throughout.

\begin{figure*}[t]
  \centering
  \includegraphics[width=\textwidth]{fig3_generalization.pdf}
  \caption{Greater-than (a): the causal sign flip replicates across scale, shaded
  95\% bootstrap intervals; the intervention helps during the dip and harms at
  maturity. Subject--verb agreement (b): accuracy dips deeply (grey, right axis)
  yet the attractor intervention stays positive at every checkpoint (band above
  zero), a boundary consistent with the shortcut not being reused. SVA uses a
  bounded probability-difference metric, so signs are compared, not magnitudes.}
  \label{fig:general}
\end{figure*}
The account also predicts a boundary: a task whose early shortcut remains
harmful, and is not later reused, should show a dip without a reversal.
Subject--verb agreement, a canonical syntactic-agreement behavior, where an
attractor noun of the opposite number intervenes between subject and verb, is
this case. Agreement accuracy dips deeply during training, to 2.5\% at step 512,
yet an analogous attractor intervention does not reverse: correcting the
attractor's number raises the agreement margin at every checkpoint, from $+0.020$
at the dip to $+0.022$ at maturity, with confidence intervals positive at every
checkpoint. Subject--verb agreement is a boundary case rather than a matched
negative control, because the task, metric, and intervention all differ; we
therefore compare signs, not magnitudes. Treating it as a boundary case rather
than overclaiming a matched control is both the more honest reading and still an
informative one: agreement is an independently motivated behavior with a deep dip
of its own, so the \emph{absence} of a reversal there is not a hand-built result.
Under the analogous attractor intervention, that absence is consistent with the
attractor not being reused later in training.

Together, greater-than and subject--verb agreement support the predicted
contrast: a sign flip when an early shortcut is later reused, and a dip without
reversal when it remains harmful. We read the subject--verb agreement result as a
boundary case rather than a matched negative control.

% ============================ 6. WHY / CONTROL (~1.3 pp) ===================
\section{Predicting and Controlling the Dip}
\label{sec:control}
If the dip reflects an imbalance between a fast copy-like signal and a slower
task-specific signal, then changing the training signal should change the dip. We
test this at two levels. The first is a two-layer toy transformer trained on a
mixture of two synthetic tasks: a \emph{copy} task, whose target is simply a
repeated token (rewarding the copy-like signal), and an \emph{IOI-like} task,
whose target requires preferring the non-repeated name (rewarding the
task-specific signal). Sweeping the fraction of copy examples, both the dip and
the sign flip appear in 9 of 10 configurations and in all three seeds, and the
copy fraction sets the timing of the dip: more copy pressure delays the point at
which the task signal overtakes it. This gives a controlled handle on the
mechanism class, though the toy model is illustrative rather than a claim about
the full circuit.

In the from-scratch setting, we test controllability of the behavioral dip, not
the full mechanistic sign flip. We train GPT-2-style models (8 layers,
width 512) from scratch on WikiText-103 with synthetic IOI examples injected at
rates of 0, 5, 15, and 30\%, evaluating on names held out from training
($n{=}3$ seeds). The injection rate predictably sets the depth of the behavioral
dip and whether the model recovers (Figure~\ref{fig:controlled}). With no
injected examples the model never recovers, sitting at a floor of 24--32\%; at 5
and 15\% it passes through a mild dip to a floor near 46\% and recovers to
roughly 90\%; at 30\% the dip disappears, with held-out accuracy near 76\%, suggesting a
generalization tradeoff rather than a monotonic benefit of more data. The floor rises
with injection rate, and the pattern replicates across seeds.

\begin{figure}[t]
  \centering
  \includegraphics[width=\columnwidth]{fig4_controlled_dip.pdf}
  \caption{From-scratch LMs: the rate of injected task examples predictably sets
  the depth of the behavioral dip and whether the model recovers (mean $\pm$ std
  over seeds); 0\% never recovers.}
  \label{fig:controlled}
\end{figure}

These from-scratch experiments do not establish the same mechanistic sign flip;
they test whether the training distribution predictably modulates the behavioral
dip, strengthening, weakening, or eliminating it. As a partial mechanistic check,
we applied the analogous S2 intervention in these models: the WikiText model,
which remains below chance throughout the run, shows positive $\Delta\mathrm{LD}$
at every checkpoint, the same sign as the pretrained dip; because these models do
not recover to high accuracy within the run, the maturity (negative) side of the
flip is not reached. The full mechanistic reversal is
demonstrated in the pretrained models
(Sections~\ref{sec:spine}--\ref{sec:mechanism}) and reproduced in the toy
setting.

% ============================ 7. RELATED WORK (~0.4 pp) ====================
\section{Related Work}
We organize related work by the explanations our results rule in or out.

\paragraph{Emergence and metric mirages.}
Apparent capability jumps with scale have been attributed to the choice of a
discontinuous metric rather than to a change in the underlying model
\citep{wei2022emergent,schaeffer2023mirage}. Our dip is the opposite situation
along the training axis: it does not disappear under a continuous metric, since
the logit difference itself goes negative. The below-chance regime is a genuine
behavioral phenomenon, not an artifact of thresholding.

\paragraph{Grokking and double descent.}
Non-monotonic generalization has been studied as grokking, where test accuracy
rises long after training loss saturates \citep{power2022grokking,
nanda2023progress}, and as double descent in the loss itself
\citep{nakkiran2021double}. In our setting the language-model loss is monotonic
throughout the dip, so the phenomenon is not loss-level non-monotonicity; it is a
transient reversal in a task behavior and in the causal role of a component,
invisible to the loss.

\paragraph{Circuit analysis across training.}
Most circuit analyses examine a model at the end of pre-training
\citep{wang2023ioi,hanna2023greater}. \citet{tigges2024circuits} track circuits
across training and scale and find the overall algorithm remarkably stable. Our
results are complementary: we characterize a transient formation window, before
that stability sets in, in which a component's causal role reverses. Both can
hold (late-training stability and a transient early reversal), and together they
argue that the formation phase deserves separate study. Our interventions are
closer to activation patching and head ablation than to full causal mediation:
path patching \citep{wang2023ioi}, automated circuit discovery
\citep{conmy2023automated}, and causal scrubbing are natural next steps for
testing whether the detected duplicate signal is routed into the later
suppression component.

\paragraph{Shortcut learning.}
Models often exploit shortcuts that correlate with the target
\citep{geirhos2020shortcut}. Our account adds a temporal dimension: a shortcut
(copying the repeated name) can be acquired early and is harmful on its own, yet
the same signal is later repurposed by a suppression component, which is why the
intervention's effect reverses rather than merely shrinking.

\paragraph{Circuits for IOI, greater-than, and agreement.}
We study the formation of behaviors whose mature circuits are known: indirect
object identification \citep{wang2023ioi}, greater-than \citep{hanna2023greater},
and subject--verb agreement, a long-standing probe of syntactic competence
\citep{linzen2016assessing,marvin2018targeted}. Induction heads provide the
prefix-matching signal we associate with detection \citep{elhage2021mathematical,
olsson2022induction}, and we analyze Pythia checkpoints \citep{biderman2023pythia}
using the unembedding-projection (logit lens) read of intermediate layers
\citep{nostalgebraist2020logitlens}.

% ============================ 8. DISCUSSION (~0.4 pp) ======================
\section{Discussion}
Our central finding is that a component's causal role is not fixed during
training: the duplicate-name representation that supports correct
behavior at maturity is one whose removal \emph{improves} behavior during the
dip. This has a direct methodological consequence. Mechanistic interpretability
that inspects a single mature checkpoint can correctly describe the final
computation yet misrepresent how it formed, including assigning a component a
role opposite to the one it held earlier. More broadly, a capability can require
sub-circuits that come online sequentially, and in the interval between them the
model can be measurably worse than chance, a failure mode that the loss curve
does not reveal. Studying behaviors as they form, rather than only once trained,
is therefore necessary for an accurate account of how language models acquire
them.

A practical implication is that circuit descriptions should be indexed by
training stage, not stated once for the final model. The same component can be
harmful, neutral, or helpful depending on whether the complementary components it
works with have yet formed, so a description that is faithful at maturity can be
exactly wrong about an earlier checkpoint. This also suggests a way to read
below-chance behavior diagnostically: a transient interval in which a model is
worse than chance can mark the gap between two sub-circuits forming, and
locating that gap points to which component is still missing. This suggests the
detection-before-suppression pattern may be one instance of a broader class in
which an early, cheaply learned signal is later gated or repurposed by a
component that forms more slowly.

% ============================ LIMITATIONS (REQUIRED, ~0.5 pp, NOT counted) =
\section*{Limitations}
We show that detection and suppression components form at different times, but
not that they are wired into a single composed circuit; demonstrating their
composition is left to future work. Our mechanistic analysis is at 160M and 410M,
while the phenomenon itself is documented at three scales. All results are on
decoder-only models, in English, on three structured tasks; we do not claim the
dip or the reversal is universal. Our IOI protocol uses curated templates and
single-token names to make the IO-versus-S contrast and the activation patching
unambiguous; extending the analysis to multi-token names and more naturalistic
contexts is an important test of ecological robustness. The toy model reproduces the dip and reversal
but is illustrative of a mechanism class rather than a faithful model of Pythia,
and the from-scratch experiments establish control over the behavioral dip, not
the same mechanistic reversal. We analyze one weights seed for 410M and 1B; for
160M, the behavioral dip reproduces across all nine PolyPythias seeds, and the
causal sign flip reproduces in an additional 160M seed; the from-scratch and toy
settings also include multiple seeds. The terms
\emph{detection} and \emph{suppression} label empirically dissociated components,
not their full internal computation. Finally, subject--verb agreement uses a
bounded probability-difference metric, so across tasks we compare the sign of the
intervention effect rather than its magnitude.

\begin{thebibliography}{}

\bibitem[Biderman et al.(2023)]{biderman2023pythia}
Stella Biderman, Hailey Schoelkopf, Quentin Anthony, Herbie Bradley, Kyle
O'Brien, Eric Hallahan, et~al. 2023.
\newblock Pythia: A suite for analyzing large language models across training
  and scaling.
\newblock In \emph{Proceedings of the 40th International Conference on Machine
  Learning (ICML)}.

\bibitem[Conmy et al.(2023)]{conmy2023automated}
Arthur Conmy, Augustine~N. Mavor-Parker, Aengus Lynch, Stefan Heimersheim, and
  Adri\`a Garriga-Alonso. 2023.
\newblock Towards automated circuit discovery for mechanistic interpretability.
\newblock In \emph{Advances in Neural Information Processing Systems 36
  (NeurIPS)}.

\bibitem[Cunningham et al.(2024)]{cunningham2024sparse}
Hoagy Cunningham, Aidan Ewart, Logan Riggs, Robert Huben, and Lee Sharkey.
  2024.
\newblock Sparse autoencoders find highly interpretable features in language
  models.
\newblock In \emph{The Twelfth International Conference on Learning
  Representations (ICLR)}.

\bibitem[Elhage et al.(2021)]{elhage2021mathematical}
Nelson Elhage, Neel Nanda, Catherine Olsson, Tom Henighan, Nicholas Joseph, Ben
  Mann, et~al. 2021.
\newblock A mathematical framework for transformer circuits.
\newblock \emph{Transformer Circuits Thread}.

\bibitem[Elhage et al.(2022)]{elhage2022superposition}
Nelson Elhage, Tristan Hume, Catherine Olsson, Nicholas Schiefer, Tom Henighan,
  Shauna Kravec, et~al. 2022.
\newblock Toy models of superposition.
\newblock \emph{Transformer Circuits Thread}.

\bibitem[Geirhos et al.(2020)]{geirhos2020shortcut}
Robert Geirhos, J\"orn-Henrik Jacobsen, Claudio Michaelis, Richard Zemel,
  Wieland Brendel, Matthias Bethge, and Felix~A. Wichmann. 2020.
\newblock Shortcut learning in deep neural networks.
\newblock \emph{Nature Machine Intelligence}, 2(11):665--673.

\bibitem[Hanna et al.(2023)]{hanna2023greater}
Michael Hanna, Ollie Liu, and Alexandre Variengien. 2023.
\newblock How does {GPT}-2 compute greater-than?: Interpreting mathematical
  abilities in a pre-trained language model.
\newblock In \emph{Advances in Neural Information Processing Systems 36
  (NeurIPS)}.

\bibitem[Linzen et al.(2016)]{linzen2016assessing}
Tal Linzen, Emmanuel Dupoux, and Yoav Goldberg. 2016.
\newblock Assessing the ability of {LSTM}s to learn syntax-sensitive
  dependencies.
\newblock \emph{Transactions of the Association for Computational Linguistics},
  4:521--535.

\bibitem[Marvin and Linzen(2018)]{marvin2018targeted}
Rebecca Marvin and Tal Linzen. 2018.
\newblock Targeted syntactic evaluation of language models.
\newblock In \emph{Proceedings of the 2018 Conference on Empirical Methods in
  Natural Language Processing (EMNLP)}.

\bibitem[Nakkiran et al.(2021)]{nakkiran2021double}
Preetum Nakkiran, Gal Kaplun, Yamini Bansal, Tristan Yang, Boaz Barak, and Ilya
  Sutskever. 2021.
\newblock Deep double descent: Where bigger models and more data hurt.
\newblock \emph{Journal of Statistical Mechanics: Theory and Experiment},
  2021(12):124003.

\bibitem[Nanda et al.(2023)]{nanda2023progress}
Neel Nanda, Lawrence Chan, Tom Lieberum, Jess Smith, and Jacob Steinhardt.
  2023.
\newblock Progress measures for grokking via mechanistic interpretability.
\newblock In \emph{The Eleventh International Conference on Learning
  Representations (ICLR)}.

\bibitem[nostalgebraist(2020)]{nostalgebraist2020logitlens}
nostalgebraist. 2020.
\newblock Interpreting {GPT}: The logit lens.
\newblock LessWrong.

\bibitem[Olsson et al.(2022)]{olsson2022induction}
Catherine Olsson, Nelson Elhage, Neel Nanda, Nicholas Joseph, Nova DasSarma,
  Tom Henighan, et~al. 2022.
\newblock In-context learning and induction heads.
\newblock \emph{Transformer Circuits Thread}.

\bibitem[Power et al.(2022)]{power2022grokking}
Alethea Power, Yuri Burda, Harri Edwards, Igor Babuschkin, and Vedant Misra.
  2022.
\newblock Grokking: Generalization beyond overfitting on small algorithmic
  datasets.
\newblock \emph{arXiv preprint arXiv:2201.02177}.

\bibitem[Schaeffer et al.(2023)]{schaeffer2023mirage}
Rylan Schaeffer, Brando Miranda, and Sanmi Koyejo. 2023.
\newblock Are emergent abilities of large language models a mirage?
\newblock In \emph{Advances in Neural Information Processing Systems 36
  (NeurIPS)}.

\bibitem[Tigges et al.(2024)]{tigges2024circuits}
Curt Tigges, Michael Hanna, Qinan Yu, and Stella Biderman. 2024.
\newblock {LLM} circuit analyses are consistent across training and scale.
\newblock In \emph{Advances in Neural Information Processing Systems 37
  (NeurIPS)}.

\bibitem[Wang et al.(2023)]{wang2023ioi}
Kevin~Ro Wang, Alexandre Variengien, Arthur Conmy, Buck Shlegeris, and Jacob
  Steinhardt. 2023.
\newblock Interpretability in the wild: a circuit for indirect object
  identification in {GPT}-2 small.
\newblock In \emph{The Eleventh International Conference on Learning
  Representations (ICLR)}.

\bibitem[Wei et al.(2022)]{wei2022emergent}
Jason Wei, Yi Tay, Rishi Bommasani, Colin Raffel, Barret Zoph, Sebastian
  Borgeaud, et~al. 2022.
\newblock Emergent abilities of large language models.
\newblock \emph{Transactions on Machine Learning Research (TMLR)}.

\end{thebibliography}

\appendix

\section{Robustness of the Sign Flip}
\label{sec:appendix}
We report three robustness checks on the Pythia-160M sign flip (dip step 2000,
mature step 143000; curated protocol, S2-to-control intervention).

\paragraph{Clustered confidence intervals.}
Because curated prompts share templates and names, example-level bootstrap
intervals can understate uncertainty. Table~\ref{tab:clustered} reports
$\Delta\mathrm{LD}$ with intervals bootstrapped three ways: over examples, over
templates (clusters), and over name pairs (clusters). All three exclude zero at
both stages and preserve the sign flip.

\begin{table}[h]
  \centering\small
  \begin{tabular}{lcc}
    \toprule
    Bootstrap clustering & $\Delta$LD (dip) & $\Delta$LD (mature) \\
    \midrule
    Example      & $+0.95$ [.86,1.05] & $-4.12$ [-4.28,-3.96] \\
    Template     & $+0.95$ [.68,1.25] & $-4.12$ [-4.54,-3.72] \\
    Name pair    & $+0.95$ [.78,1.15] & $-4.12$ [-4.45,-3.82] \\
    \bottomrule
  \end{tabular}
  \caption{Clustered 95\% bootstrap intervals; the flip survives all three.}
  \label{tab:clustered}
\end{table}

\paragraph{Layer-window stability.}
The reported window $\{3,4,5\}$ is not special. Table~\ref{tab:windows} sweeps
neighboring windows; every window is positive at the dip and negative at
maturity. Magnitude decays for deeper windows, but the sign flip does not depend
on the window.

\begin{table}[h]
  \centering\small
  \begin{tabular}{lcc}
    \toprule
    Layer window & $\Delta$LD (dip) & $\Delta$LD (mature) \\
    \midrule
    $\{2,3,4\}$   & $+0.95$ & $-4.11$ \\
    $\{3,4,5\}$   & $+0.95$ & $-4.12$ \\
    $\{4,5,6\}$   & $+0.80$ & $-4.12$ \\
    $\{5,6,7\}$   & $+0.73$ & $-3.93$ \\
    $\{3,4,5,6\}$ & $+0.95$ & $-4.12$ \\
    \bottomrule
  \end{tabular}
  \caption{Sign flip is stable across neighboring layer windows.}
  \label{tab:windows}
\end{table}

\paragraph{Selection-bias-robust trajectory.}
Identifying the dip checkpoint on a selection split of names and reporting on a
disjoint held-out split, the held-out $\Delta\mathrm{LD}$ is
$+0.003$ (step 0), $+0.028$ (512), $+0.219$ (1000),
$\mathbf{+1.13}$ (2000, dip), and $\mathbf{-4.09}$ (143000, mature), confirming
the flip on names never used for checkpoint selection.

\section{Generalization Details}
\paragraph{Greater-than across scale.}
The greater-than intervention effect (start-year representation patched toward a
control) is positive during the dip and negative at maturity at all scales:
$+0.08\to-0.74$ (160M), $+0.10\to-0.67$ (410M), $+0.32\to-0.73$ (1B); all
intervals exclude zero. As a negative control, next-token loss on natural text is
monotonic across the same checkpoints.

\paragraph{Subject--verb agreement trajectory.}
Table~\ref{tab:sva} gives the full subject--verb agreement trajectory: accuracy
dips to 2.5\% yet the attractor intervention $\Delta$ stays positive throughout,
with intervals excluding zero, so the role does not reverse. The metric is a
bounded probability difference, so we compare signs, not magnitudes, with the
logit-difference results elsewhere.

\begin{table}[h]
  \centering\small
  \begin{tabular}{lcc}
    \toprule
    Step & SVA accuracy & $\Delta$ (intervention) \\
    \midrule
    256    & 12.5\% & $+0.020$ \\
    512    & \phantom{0}2.5\% & $+0.010$ \\
    1000   & 42.5\% & $+0.015$ \\
    2000   & 68.0\% & $+0.014$ \\
    5000   & 95.0\% & $+0.014$ \\
    13000  & 96.5\% & $+0.018$ \\
    143000 & 88.5\% & $+0.022$ \\
    \bottomrule
  \end{tabular}
  \caption{Subject--verb agreement: deep dip, no reversal.}
  \label{tab:sva}
\end{table}

\section{Mechanism Details}
\paragraph{Suppression heads.}
We rank attention heads by mean-ablation effect on mature LD. The five strongest
inhibitory heads are L8H9 ($-2.66$), L3H0 ($-1.78$), L4H11 ($-1.54$), L6H2
($-0.78$), and L1H8 ($-0.66$); we track L8H9, the strongest, across training.
Its ablation effect is near zero at the dip and reaches $-2.66$ at maturity,
approximately 70\% of the mature base LD magnitude of $+3.79$.

\paragraph{Detection signal and logit lens.}
The prefix-matching (induction) score rises from $0.04$ to $0.94$ by step 1000
and remains high through the below-chance interval. An unembedding projection of
intermediate layers favors S at every layer during the dip; a late layer (L8)
first favors IO at recovery. The layer marker is clean at 160M and noisier at
410M. Circuit-discovery and projection methods follow standard practice
\citep{conmy2023automated,nostalgebraist2020logitlens}.

\paragraph{Composition (preliminary path patching).}
Path patching at maturity indicates that the dominant suppression head L8H9 does
not act in isolation: ablating it changes the contribution of downstream
late-layer heads, shifting probability from IO toward S2 (for L9H1,
$\Delta\mathrm{IO}=-0.17$, $\Delta\mathrm{S2}=+0.17$; similarly for L9H7), and
early heads (L0H0, L3H0) feed downstream name-mover heads (L10H5, L9H4). This is
consistent with a suppression-to-name-mover route. We report it as preliminary:
establishing the full detection-to-suppression path would require path patching
across the dip, which we leave to future work.

\section{Negative Controls and Alternatives}
\paragraph{Subspace recovery.}
Projecting onto ten directions at layer~4 recovers approximately 107\% of the
dip's logit-difference deficit, indicating the dip signal is low-rank and
linearly accessible. The same directions estimated on one task do not transfer to
another (cross-task projection fails), and a sparse-autoencoder feature analysis
yields no single feature accounting for the reversal
\citep{elhage2022superposition,cunningham2024sparse}; we report these as honest
negative results.

\paragraph{Grokking and positional controls.}
Across nine grokking runs, the ablation effect is never positive during the
generalization transition, so the sign flip is not a generic grokking signature.
As a positional control, applying the duplicate-name patch at a non-duplicate
position leaves the logit difference essentially unchanged ($\Delta \approx 0$
at the dip), confirming the effect is specific to the S2 position rather than a
generic consequence of perturbing the residual stream.

\section{Additional Models}
\label{sec:appendix-models}
Both the below-chance dip and the causal sign flip reproduce beyond the three
main Pythia runs.

\paragraph{Nine PolyPythias seeds.}
We evaluate all nine released PolyPythias-160M seeds (three random seeds, three
data-order seeds, three weight-initialization seeds). Every one dips below chance
at step~2000, with floor accuracy ranging from 18.0\% to 36.7\% and a negative
mean logit difference in all nine; recovery to 95--100\% follows in every case
(Figure~\ref{fig:polyseeds}). The dip is therefore not an artifact of the
particular seed used for the main results.

\begin{figure}[t]
  \centering
  \includegraphics[width=\columnwidth]{fig5_polypythias_seeds.pdf}
  \caption{All nine PolyPythias-160M seeds (light; mean in dark) dip below the
  50\% forced-choice baseline near step~2000 before recovering, confirming the
  dip is not seed-specific.}
  \label{fig:polyseeds}
\end{figure}

\paragraph{Sign flip across seed and architecture.}
Table~\ref{tab:crossmodel} reports the causal intervention for a different
Pythia-160M random seed (PolyPythias seed-1) and an independently trained,
different-codebase GPT-2 (Stanford CRFM \texttt{alias-gpt2-small-x21}): both dip
below chance and both show the same $\Delta\mathrm{LD}$ reversal, positive at the
floor and negative at maturity.

\begin{table}[h]
  \centering\small
  \setlength{\tabcolsep}{4pt}
  \begin{tabular}{lccc}
    \toprule
    Model & Floor & $\Delta$LD & $\Delta$LD \\
          & acc. & (floor) & (mature) \\
    \midrule
    PolyPythias seed-1 & 30.7\% & $+0.55$ & $-3.92$ \\
    Stanford GPT-2     & 12.0\% & $+1.03$ & $-2.89$ \\
    \bottomrule
  \end{tabular}
  \caption{The dip and the sign flip replicate on a different random seed and a
  different architecture/codebase. Both show $\Delta\mathrm{LD}>0$ at the dip
  floor and $\Delta\mathrm{LD}<0$ at maturity, matching the main Pythia results.}
  \label{tab:crossmodel}
\end{table}

\paragraph{A different model family.}
The dip is not confined to the Pythia family. In OLMo-1B, IOI accuracy is 44\%
with a negative mean logit difference at 4B training tokens and recovers to
61--66\% at later checkpoints, the same below-chance-then-recover shape. We did
not run the causal intervention on OLMo, so we report this as behavioral evidence
of the dip in a separate family, not of the sign flip.

\paragraph{From-scratch and other architectures.}
The dip further appears in from-scratch GPT-2-style models, where WikiText-103
training produces a deep dip without recovery and OpenWebText produces a milder
dip with partial recovery, a data-diversity gradient. GPT-2-small, GPT-2-medium,
GPT-Neo-125M, and OPT-125M all perform IOI at maturity.

\section{Probe Decomposition and Position Analysis}
\label{sec:appendix-probe}
We distinguish two probes. A \emph{duplicate-detection} probe predicts whether the prompt contains a repeated name, read from the final-token representation;
trained with cross-validation on a balanced set of prompts with and without a
repeated name (chance 50\%), its held-out
accuracy is already high during the dip (step~2000: 0.52, 0.60, 0.63, 0.72, 0.77
for layers 0--4, then 0.92, 0.92, 0.91, 0.91, 0.91, 0.91, 0.93 for layers 5--11),
and reaches ceiling by step~5000. An \emph{answer} probe, predicting which name
is the indirect object, instead achieves high in-distribution accuracy but drops
to 44--51\% on held-out names, indicating that number reflects memorization;
the causal sign flip, by contrast, is reported throughout on held-out names. The
contrast is consistent with duplicate detection being a learned structural
feature that generalizes, while the specific answer mapping is partly memorized.
A position scan shows the S1 position carries more decision-relevant signal than
S2 at maturity, and the dominant name-mover and S-inhibition heads differ between
the original and retrained checkpoints, consistent with component-level
fluctuation under a stable algorithm.

\section{Reproducibility and Compute}
All experiments use released Pythia checkpoints
\citep{biderman2023pythia} and the from-scratch models described in
Section~\ref{sec:control}. Pretrained-model analyses run on a single A100 (40GB);
each from-scratch run trains an 8-layer, width-512 model for 40K steps. Code,
result files, generated checkpoints, and scripts to regenerate every figure and
table are available at an anonymized repository:
\url{https://anonymous.4open.science/r/ioi-sign-flip-6FD9}.

\end{document}
